% Part:sets-functions-relations
% Chapter: size-of-sets
% Section: zig-zag

\documentclass[../../../include/open-logic-section]{subfiles}

\begin{document}

\olfileid{sfr}{siz}{zigzag}

\olsection{د کانتور د زيګزاګ طريقه}

\begin{explain}
مخکښې مو ځينې «اسانې» شمېرنې په پام کښې نيولې دي۔ اوس به
يو لږ ګران کار په پام کښې ونيسو۔ د طبيعي عددونو د جوړو سټ
په پام کښې ونيسئ\oliflabeldef{sfr:set:pai:sec}{، چې په
\olref[set][pai]{sec} کښې مو داسې تعريف کړے و:}{چې داسې تعريف شوے دے:}
\[
\Nat \times \Nat = \Setabs{\tuple{n,m}}{n,m \in \Nat}
\]
دا مرتبې جوړې په يوه \emph{جدولي ترتيب} کښې داسې برابرولے شو:
\[
\begin{array}{ c | c | c | c | c | c}
& \mathbf 0 & \mathbf 1 & \mathbf 2 & \mathbf 3 & \dots \\
\hline
\mathbf 0 & \tuple{0,0} & \tuple{0,1} & \tuple{0,2} & \tuple{0,3} & \dots \\
\hline
\mathbf 1 & \tuple{1,0} & \tuple{1,1} & \tuple{1,2} & \tuple{1,3} & \dots \\
\hline
\mathbf 2 & \tuple{2,0} & \tuple{2,1} & \tuple{2,2} & \tuple{2,3} & \dots \\
\hline
\mathbf 3 & \tuple{3,0} & \tuple{3,1} & \tuple{3,2} & \tuple{3,3} & \dots \\
\hline
\vdots & \vdots & \vdots & \vdots & \vdots & \ddots\\
\end{array}
\]
په ښکاره، په $\Nat \times \Nat$ کښې هره مرتبه جوړه به په
دې جدول کښې په دقيقه توګه يو ځل راشي۔ په ځانګړي ډول،
$\tuple{n,m}$ به په $n$م قطار او $m$م ستون کښې راشي۔ خو د
داسې جدول غړي په يوه «يو بُعدي» لړ کښې څنګه برابر کړو؟ په
لاندې جدول کښې نمونه د دې کار يوه لار ښيي، که څه هم البته
ډېرې نورې لارې هم شته:
\[
\begin{array}{ c | c | c | c | c | c | c}
& \mathbf 0 & \mathbf 1 & \mathbf 2 & \mathbf 3 & \mathbf 4 &\dots \\
\hline
\mathbf 0 & 0  & 1& 3 & 6& 10 &\ldots \\
\hline
\mathbf 1 &2 & 4& 7 & 11 & \dots &\ldots \\
\hline
\mathbf 2 & 5 & 8 & 12 & \ldots & \dots&\ldots \\
\hline
\mathbf 3 & 9 & 13 & \ldots & \ldots & \dots & \ldots \\
\hline
\mathbf 4 & 14 & \ldots & \ldots & \ldots & \dots & \ldots \\
\hline
\vdots & \vdots & \vdots & \vdots & \vdots&\ldots & \ddots\\
\end{array}
\]\noindent
دا نمونه \emph{د کانتور د زيګزاګ طريقه} بلل کېږي۔ دا د
$\Nat \times \Nat$ شمېرنه داسې کوي:
\[
\tuple{0,0}, \tuple{0,1}, \tuple{1,0}, \tuple{0,2}, \tuple{1,1},
\tuple{2,0}, \tuple{0,3}, \tuple{1,2}, \tuple{2,1}, \tuple{3,0}, \dots
\]
او له دې لاندې خبره ثابتېږي:
\end{explain}

\begin{prop}\ollabel{natsquaredenumerable}
$\Nat \times \Nat$ يو !!{enumerable} سټ دے۔
\end{prop}

\begin{proof}
فرض کړئ $f \colon \Nat \to \Nat\times\Nat$ هر $k \in \Nat$
له هغې مرتبې څوګونې $\tuple{n,m} \in \Nat \times \Nat$ سره
تړي چې $k$ د کانتور د زيګزاګ په جدول کښې د $n$م قطار او
$m$م ستون قيمت وي۔
\end{proof}

\begin{explain}
دا تخنيک په ښه ډول عمومي کېدے هم شي۔ د بېلګې په توګه، د
طبيعي عددونو د مرتبو درې ګونو د سټ د شمېرنې دپاره ئې
کارولے شو؛ يعنې:
\[
\Nat \times \Nat \times \Nat = \Setabs{\tuple{n,m,k}}{n,m,k \in \Nat}
\]
$\Nat \times \Nat \times \Nat$ له $\Nat$ سره د $\Nat \times \Nat$
د کارتيزي حاصل ضرب په توګه ګڼو؛ يعنې:
\[
\Nat^3 = (\Nat \times \Nat) \times \Nat =
\Setabs{\tuple{\tuple{n,m},k}}{n, m, k
  \in \Nat }
\]
نو د يوه محور د $\Nat$ په شمېرنې او د بل محور د $\Nat^2$
په شمېرنې په نښه کولو سره، د يوه جدول په وسيله د $\Nat^3$
شمېرنه کولے شو:
\[
\begin{array}{ c | c | c | c | c | c}
& \mathbf 0 & \mathbf 1 & \mathbf 2 & \mathbf 3 & \dots \\
\hline
\mathbf{\tuple{0,0}} & \tuple{0,0,0} & \tuple{0,0,1} & \tuple{0,0,2} & \tuple{0,0,3} & \dots \\
\hline
\mathbf{\tuple{0,1}} & \tuple{0,1,0} & \tuple{0,1,1} & \tuple{0,1,2} & \tuple{0,1,3} & \dots \\
\hline
\mathbf{\tuple{1,0}} & \tuple{1,0,0} & \tuple{1,0,1} & \tuple{1,0,2} & \tuple{1,0,3} & \dots \\
\hline
\mathbf{\tuple{0,2}} & \tuple{0,2,0} & \tuple{0,2,1} & \tuple{0,2,2} & \tuple{0,2,3} & \dots\\
\hline
\vdots & \vdots & \vdots & \vdots & \vdots & \ddots \\
\end{array}
\]
نو د کانتور د زيګزاګ طريقې ته ورته طريقې په کارولو سره د~$\Nat^3$
شمېرنه هم ترلاسه کولے شو۔ او همداسې مخ په وړاندې تللے شو،
چې د هر طبيعي عدد $n$ دپاره د $\Nat^n$ شمېرنې ترلاسه کړو۔ نو لرو:
\end{explain}

\begin{prop}
د هر $n \in \Nat$ دپاره، $\Nat^n$ يو !!{enumerable} سټ دے۔
\end{prop}

\begin{prob}\label{sfr:siz:zigzag:prob:posint-n}
وښيئ چې د هر $n \in \Nat$ دپاره، $(\PosInt)^n$ يو !!{enumerable} سټ دے۔
\end{prob}

\begin{prob}\label{sfr:siz:zigzag:prob:posint-star}
  وښيئ چې $(\PosInt)^*$ يو !!{enumerable} سټ دے۔
  \cref{sfr:siz:zigzag:prob:posint-n} فرضولے شئ۔
\end{prob}

\end{document}
